% ==============================================================================
% PAGE 11: CHAPTER-1 INTRODUCTION (Numbered Page 3)
% ==============================================================================
\begin{center}
  {\fontsize{15}{18}\selectfont \bfseries \textcolor{red!80!black}{\underline{CHAPTER-1}}}\\[2mm]
  {\fontsize{16}{20}\selectfont \bfseries \textcolor{red!80!black}{\underline{INTRODUCTION}}}\\[6mm]
\end{center}

\subsection*{1.1 OVERVIEW}
\noindent
\begin{spacing}{1.2}
{\fontsize{10.5}{14.5}\selectfont
Energy plays a vital role in the development of any country. The increasing demand for energy and the depletion of conventional energy sources like coal, oil, and natural gas have made it necessary to explore renewable energy resources. Among these, solar energy is the most abundant, clean, and sustainable form of energy available.\\[3mm]
The Sun radiates a huge amount of energy in the form of light and heat. If effectively utilized, this energy can fulfill a major part of the world's energy needs. Solar energy can be converted into useful heat energy or electrical energy through various technologies such as solar water heaters, photovoltaic cells, and solar cookers.\\[3mm]
The present project, titled \textbf{``Design and Fabrication of Parabolic Water Heater with Solar Tracking and Power Generation,''} is a step toward utilizing this renewable energy for practical use. The system focuses on capturing solar radiation using a parabolic collector to heat water efficiently and uses solar tracking to improve performance. It also integrates a small-scale power generation unit, making it a multipurpose and eco-friendly system.
}
\end{spacing}

\vspace{5mm}
\subsection*{1.2 NEED OF PROJECT}
\noindent
\begin{spacing}{1.2}
{\fontsize{10.5}{14.5}\selectfont
In most rural and semi-urban areas, people still depend on firewood, kerosene, and LPG gas for heating water. These fuels not only increase household expenses but also contribute to pollution and environmental degradation.\\[3mm]
Using solar energy for heating water is an ideal solution because:
\begin{itemize}[leftmargin=8mm, itemsep=1.5mm, label=\textbullet]
  \item It is freely available and inexhaustible.
  \item It reduces dependency on non-renewable energy sources.
  \item It helps save fuel and electricity.
  \item It lowers carbon emissions and environmental pollution.
\end{itemize}
However, most existing solar water heaters are fixed in one direction and cannot follow the Sun's movement throughout the day. This results in a considerable loss of efficiency. To overcome this problem, a solar tracking mechanism is introduced to maintain the collector's orientation toward the Sun, thereby increasing heat gain and system output.
}
\end{spacing}
\clearpage
